\documentclass{article}
\usepackage[a4paper, total={6in,8in}]{geometry}
\usepackage{graphicx,float,hyperref}
\usepackage{outlines,enumerate}
\usepackage{amsmath,amssymb,amsthm}

\renewcommand{\vec}[1]{%
	\mathbf{#1}%
}

\hypersetup{
	colorlinks=true,
	linkcolor=blue,
	urlcolor=blue
}

\begin{document}
	\href{https://mycourses.unh.edu/courses/152824/files/folder/lecture?preview=15941715}{\textsf{Reference lecture (click)}}\\
	\begin{center}
		{\huge\bf Gradient Descent}
	\end{center}
	\section{Optimization}
	\begin{outline}
		\1 \textbf{Objective} is the function we are optimizing: $\mathcal{O}(\vec{w})$
		\1 Optimizing can mean either minimizing or maximizing. Optimization literature focuses on minimizing. With focus on ML we may want to:
		\2 maximize accuracy $\mathcal{O}(...) = \frac{\text{correct predictions}}{\text{all}}$
		\2 maximize likelihood $\mathcal{O}(...) = \mathcal{L}(...)$
		\1 Phrase your "wish" in the form of the objective $\mathcal{O}$
		\1 It is the job of the optimization algorithm to make your wish come true!
	\end{outline}
	\subsection{Optimizing: Maximizing vs Minimizing}
	Tricks
	\begin{outline}
		\1 $\log \mathcal{L}(\vec{w})$ and $\mathcal{L}(w)$ have max at the same position
		\1 maximize $\mathcal{L}$ or minimize $-\mathcal{L}(\vec{w})$ (negative)
		\1 $\mathcal{O}(\vec{w}) = \underbrace{-\log \mathcal{L}_{\mathcal{D}}(\vec{w})}_{\text{negative log likelihood}}$
	\end{outline}
	\subsection{Optimizing: Sloppy Notations}
	\begin{outline}
		\1 The likelihood $\mathcal{L}$ depends on data $\mathcal{D}$ such as $\vec{X}$ and $\vec{y}$ and/or hyperparameters $\alpha, \beta$
		\1 We gloss over constants like $\mathcal{D}$ when we define the objective function
		\[ \mathcal{O}(\vec{w}) = -\log\mathcal{L}_{\mathcal{D}}(\vec{w}) = -\log \mathcal{L}(\vec{y} \;|\; \vec{X}, \vec{w}) \]
	\end{outline}
	These all mean the same thing.\\
	We might also say "minimize $-\log \mathcal{O}(\vec{w})$"
	\section{Partial Derivative and Gradient}
	\begin{outline}
		\1 Given a multi-variate function such as $f(x,y)$ or $f(\vec{w}$ or $\mathcal{O}(\vec{w})$
		\1 The \textbf{partial derivative} $\frac{\delta f(x,y)}{\delta x}$ is the normal derivative of the function $f(x)$, where $y$ is interpreted as a constant.
	\end{outline}
	Example:
	\begin{align*}
		f(x,y) = 3 + x^3y^2 + y^2 + 7x\\
		\frac{\delta f(x,y)}{\delta x} = 3y^2x^2 + 7\\
		\frac{\delta f(x,y)}{\delta y} = 2x^3y + 2y
	\end{align*}
	The \textbf{gradient} of the function is the vector of partial derivatives
	\[ \nabla f(x,y) \begin{bmatrix} \frac{\delta f(x,y)}{\delta x} \\ \frac{\delta f(x,y)}{\delta y} \end{bmatrix} =
	\begin{bmatrix} 3y^2x^2 + 7\\2x^3y + 2y \end{bmatrix} \]
	We can evaluate the gradient at any point $x,y$.
	\subsection{Gradient Descent}
	Goal: $\arg \min_{\vec{w}} f_{\mathcal{D}}(\vec{w})$ - function also includes $\vec{x}_i, y_i$ as constants
	\begin{outline}
		\1 Choose a random point $\vec{w}^{(0)} \in \mathbb{R}^J$
		\1 For $t = 1$ to ... until converged:
		\2 Identify upwards direction at current point: $\vec{\tilde{d}} = \nabla f(\vec{w}^{(t-1)})$
		\2 Take a tiny step (of length $\alpha$) in opposite direction\\
		$\vec{w}^{(t)} = \vec{w}^{(t-1)} - \alpha \vec{\tilde{d}}$
	\end{outline}
	\section{Hyperplane Models Revisited}
	\subsection{Linear Model Objective}
	\begin{outline}
		\1 Given training data $\mathcal{D} = \{(\vec{x}_i, y_i)\}_{i=1}^N$
		\1 Find a weight vector $\vec{w}$ that minimizes the errors
		\1 Predict $i$ as positive if $\vec{w}\vec{x}_i \geq 0$
		\1 Since only signs matter, this can be written compactly
		\[ \min_{\vec{w}} \sum_i \Bigg\| \underbrace{\underbrace{y_i}_{\text{truth}} \cdot \underbrace{(\vec{w}\vec{x}_i)}_{\text{predict}} < 0}_{\text{true=mistake?}} \Bigg\| \]
		if we define target variable $y_i = \begin{cases} +1 & \text{positive}\\
		-1 & \text{negative} \end{cases}$
		\1 Indicator function: $\|$True$\| = 1$ vs $\|$False$\| = 0$
		\2 Actual $y_i = +1$, prediction $\vec{w}\vec{x}_i = 0.5$ then $+1 \cdot 0.5 \nless 0$
		\2 Actual $y_i = -1$, prediction $\vec{w}\vec{x}_i = 0.5$ then $-1 \cdot 0.5 < 0$
		\2 Actual $y_i = -1$ prediction $\vec{w}\vec{x}_i = -0.5$ then $-1 \cdot (-0.5) \nless 0$
	\end{outline}
	\subsection{Zero-One Loss}
	\begin{outline}
		\1 In Optimization, minimize objective function
		\1 Minimizing the error rate
		\[ \min_\vec{w} \underbrace{\Bigg\| y_i(\vec{w}\vec{x}_i) < 0}_{\text{Objective } \mathcal{O}} \]
		\1 Also called the \textbf{Zero-One Loss}, awarding 0 or 1 per example
		\1 We want as many zeros as possible, because 1s indicate mistakes. Remember: we are minimizing Objective $\mathcal{O}$!
	\end{outline}
\end{document}